\section{Discussion}\label{sec:discussion}

\textbf{Distance to real time decoding.} A superconducting device produces a syndrome round about every microsecond
\cite{google2025below}. To keep up with one logical qubit, a pre-decoder that consumes windows of $T=d$ rounds must
finish at least one window every $d$ microseconds: \NumRealTimeWpsFast{} windows per second for the Fast network at
$d=9$ and \NumRealTimeWpsAcc{} for the Accurate network at $d=13$. Because the work per window is fixed, this asks for
\NumRealTimeTmacFast{}$\times10^{12}$ and \NumRealTimeTmacAcc{}$\times10^{12}$ multiply accumulate operations per
second. Four cards deliver \NumTmacFast{}$\times10^{12}$ and \NumTmacAcc{}$\times10^{12}$, and even if every one of the
9{,}024 DSP blocks of a U55C \cite{amd_ds978} performed two 8 bit products per cycle at 300~MHz, one card would reach
\NumDevicePeakTmac{}$\times10^{12}$. The released networks evaluated densely on full windows therefore need more than ten
times the peak arithmetic of a card of this class for real time use, independently of how well the engine is
scheduled; our kernel already reaches \NumModelEffFast{}\% and \NumModelEffAcc{}\% of the rate its schedule allows. The
measured latency of \NumLatFast{}~ms per window leads to the same conclusion. Real time pre-decoding in hardware
therefore needs networks designed for it: fewer channels, fewer layers, or sliding windows that reuse the results of
earlier rounds instead of recomputing a full window, and possibly the sparsity of the syndrome, which in our kernel
only the first layer could exploit because the hidden activations after GELU are dense. On platforms with millisecond
syndrome cycles, a CPU already decodes in real time \cite{ye2026megaquop}, and a pre-decoder is not needed for speed.

\textbf{What the cards are useful for.} The design is a throughput engine. Evaluating a pre-decoder on many shots, as
in Table~\ref{tab:ler}, and generating pre-decoded syndromes for decoder studies are batch tasks of independent windows,
which spread over the cards without communication and scale linearly with their number. For these tasks, however,
the advantage of the cards over the host is small: four cards are \NumSpeedFast{} (Fast) and \NumSpeedAcc{}
(Accurate) times faster than PyTorch on the two socket, 128 core host at its best thread count and batch size, so the
host matches \NumHostCardsFast{} cards for the Fast network and \NumHostCardsAcc{} for the Accurate network, and for a single window the CPU is as fast (Fast) or faster (Accurate). What the cards
offer is that the work leaves the host, whose cores then remain free for the matching decoder, and an integer output
known to be bit exact. The CPU baseline is the floating point network in PyTorch, the form in which the networks are
released; an optimized integer CPU implementation could close or reverse the remaining gap, and we did not measure GPUs
or the power of either platform. The per card throughput is limited by the single engine of 1{,}024 products per cycle
per instance, which uses less than half of the DSP blocks of an SLR (Table~\ref{tab:resources}); a larger engine or
two products per DSP block would raise it, but not by the factor that real time use needs.

\textbf{Accuracy of the integer model.} The integer model keeps the gain of the Fast network and about two thirds of
that of the Accurate network: the Fast network at $p=0.006$, $d=9$ gives \NumRatioFastIntNineHigh{} against \NumRatioFastFpNineHigh{}, and the Accurate
network at $p=0.006$, $d=13$ gives \NumRatioAccIntThirteenHigh{} against \NumRatioAccFpThirteenHigh{}, both over both
bases. Calibration uses windows of both bases at the noise level of the evaluation and no retraining; quantization aware
training could recover the remainder.

\textbf{Limits of this study.} The kernel holds windows up to $13\times13\times13$ sites and the two released surface
code networks; larger windows need tiling of the activation buffers. Every launch reloads the weights from HBM, which
costs little in throughput but adds to the latency of a single window; a kernel that keeps the weights resident between
launches would remove that part. Logical error rates are measured for memory experiments under NVIDIA's circuit level
noise model at two noise levels, with the number of shots given in Table~\ref{tab:ler}; the smallest counts, at
$p=0.003$, carry a relative statistical uncertainty of up to about ten percent.
